\chapter*{INTRODUCTION GÉNÉRALE}
\markboth{\MakeUppercase{INTRODUCTION GÉNÉRALE}}{}
%\addstarredchapter{INTRODUCTION GÉNÉRALE}
\addcontentsline{toc}{chapter}{INTRODUCTION GÉNÉRALE}
\adjustmtc
\thispagestyle{MyStyle}

\indent Aujourd’hui, plusieurs entreprises de différentes spécialités sont prêtes à investir des sommes considérables dans l’implantation de nouvelles technologies logicielles. Ces entreprises ont pour finalité d’améliorer leurs services, d’accroitre leur agilité et flexibilité, de réduire leurs coûts, d’augmenter leur production ainsi que de faire face aux défis du marché. Cependant, la diversification et la croissance des activités au sein des entreprises ont fait grandir le besoin d’intégrer plusieurs technologies. La tâche de gérer efficacement toutes ces activités s’avère ainsi de plus en plus complexe et difficile. Pour surmonter ses propres difficultés, une entreprise a besoin d’utiliser des outils optimisés qui s’adaptent à la réalisation de toutes les activités en offrant des fonctionnalités riches et utiles. Parmi ces outils, nous retrouvons les systèmes intégrés de gestion tels que les ERP (Entreprise Ressources Planning) Les ERP sont des outils de gestion et d’analyse permettant d’optimiser la diffusion de l’information en interne, d’améliorer les processus de gestion et d’automatiser les activités et les tâches de l’entreprise.\par

Il en va de même pour l’entreprise \reportEntreprise de plasturgie qui souhaite optimiser sa gestion de production par un système d'information unique grâce à un progiciel de gestion intégré L'intérêt de notre projet est d'automatiser la gestion de la production à l'aide de l'ERP. Afin de développer ce projet, il est nécessaire de passer par une étape d'analyse des besoins, suivie d'une conception détaillée du projet, puis de se concentrer sur le développement de modules, en sachant répondre à tous exigence.\par
Pour la société \reportEntreprise , le déploiement d'un ERP est un levier de digitalisation indispensable\\ pour moderniser les processus et se conformer aux standards nationaux et internationaux.\\
C’était une mission qui semblait importante pour la société, ce qui nous a poussés à nous investir énormément pour la mener à bien.\\
Le présent rapport est organisé en quatre chapitres : 
le premier chapitre est un cadrage générale du projet qui présente un aperçu sur les différents concepts théoriques d’un système ERP, objectifs et enjeux spécifique à la société \reportEntreprise.\\
Le second chapitre est une analyse et spécifications des besoins : problématique, identifications des acteurs et détail des besoins de l’application.\\
Dans le troisième chapitre, nous allons présenter une conception de l’application qui comprend l’architecture technique du système, choix de conception et modèles utilisées pour structurer l’application Ce travail établit la feuille de route logicielle. En fixant les règles et la structure en amont, il facilite et fiabilise l'exécution technique du projet.\\
Finalement, dans le quatrième chapitre nous dresserons le bilan de ce modeste travail dans la conclusion de ce rapport.


